% ============================================================
%  Objetivo 3 — PROSITE y patrones
% ============================================================

% ============================================================
\section{Ejercicio 3a}

\begin{consigna}
¿Qué tipo de signature es PS50111? ¿A qué dominio corresponde?
\end{consigna}

PS50111 es un \textbf{perfil} (\textit{profile}), es decir una matriz de pesos posición-específica
que captura cuánto varía cada posición del dominio --- no un patrón de expresiones regulares. Los
perfiles son más sensibles que los patrones porque modelan variabilidad posición a posición en lugar
de buscar residuos fijos.

El resultado de ScanProsite encontró \textbf{1 hit}:

\begin{table}[H]
\centering
\footnotesize
\begin{tabular}{llP{2.5cm}P{5cm}}
\toprule
Signature & Posiciones & Score & Dominio predicho \\
\midrule
PS50111 & 331--588 & 29.836 & ``Methyl-accepting transducer'' (MCPsignal) \\
\bottomrule
\end{tabular}
\end{table}

PS50111 corresponde al dominio \textbf{MCPsignal}, no al HNOB. Las posiciones 331--588 coinciden
exactamente con el dominio C-terminal de señalización de quimiotaxis. Algo a notar: el dominio
HNOB (posiciones 1--190) \textbf{no fue detectado} por PROSITE. Pfam e InterPro sí lo detectaron,
pero PROSITE no tiene un perfil de HNOB que alcance el umbral para esta secuencia. Esto muestra
que distintas bases de datos pueden cubrir distintos dominios, y que usar una sola herramienta
no siempre es suficiente.

% ============================================================
\section{Ejercicio 3b}

\begin{consigna}
Construya a mano un patrón a partir del logo del HMM. Búsquelo en PROSITE y describa qué tipo
de proteínas encuentra.
\end{consigna}

Construimos el patrón a partir de las tres posiciones más conservadas del HMM de HNOB (His
proximal --- pos.\,103, Pro --- pos.\,116, Tyr --- pos.\,132):

\begin{center}
\ttfamily H-x(11,13)-P-x(14,16)-Y
\end{center}

Al buscarlo en Swiss-Prot obtuvimos \textbf{más de 10.000 hits} (el servidor alcanzó el límite).
El número esperado de matches al azar es $\sim$14.021 en 100.000 secuencias, lo que indica que el
patrón es prácticamente no específico.

Entre los resultados aparecen proteínas completamente no relacionadas: virus de insectos (IIV-3,
Frog virus), proteínas de cerdo africano (ASFV), lectinas, glucosidasas, enzimas de reparación de
ADN de docenas de organismos distintos. No encontramos ninguna ``methyl-accepting chemotaxis
protein'' ni proteína de unión a NO entre los primeros resultados.

El problema es que el patrón tiene solo 3 posiciones fijas (H, P, Y) separadas por ventanas muy
largas ($\sim$30 residuos cada una), lo que permite demasiados matches al azar. Esto ilustra bien
la diferencia entre un patrón manual y un perfil HMM: el perfil usa la información de todas las
posiciones del alineamiento en simultáneo, no solo 3 residuos ancla.

% ============================================================
\section{Ejercicio 3c}

\begin{consigna}
Corra PRATT con el alineamiento seed de HNOB. Compare los patrones antes y después del
refinamiento. ¿Puede localizar el patrón en el logo del HMM?
\end{consigna}

Para correr PRATT necesitamos el alineamiento seed de HNOB en formato FASTA. Descargamos el
seed desde Pfam en formato Stockholm y lo convertimos con el script
\texttt{scripts/stockholm\_to\_fasta.py}:

\begin{lstlisting}[language=bash]
python scripts/stockholm_to_fasta.py data/HNOB_seed.sto resultados/HNOB_seed.fasta
\end{lstlisting}

El seed original tenía 121 secuencias, pero PRATT acepta como máximo 100, así que tomamos
las primeras 100 y las guardamos en \texttt{resultados/HNOB\_seed\_100.fasta}. Corrimos PRATT
con Min Percentage = 70 y Pattern Refinement = ON.

\textbf{Antes del refinamiento}, los patrones con mayor fitness son muy cortos y con
posiciones degeneradas:

\begin{table}[H]
\centering
\footnotesize
\begin{tabular}{P{5.5cm}P{2cm}P{2cm}}
\toprule
Patrón & Fitness & Secuencias \\
\midrule
\texttt{Y-x-S-x-R}        & 12.5102 & 93 \\
\texttt{Y-x-S-x(1,2)-R}   & 12.0102 & 93 \\
\texttt{F-x(4,5)-D-x(1,2)-H} & 11.5102 & 76 \\
\bottomrule
\end{tabular}
\caption{Patrones de PRATT antes del refinamiento (Top 3 por fitness).}
\end{table}

El patrón más frecuente antes del refinamiento (\texttt{Y-x-S-x-R}) captura 93 de las 100
secuencias pero tiene solo 3 posiciones fijas separadas por wildcards de longitud fija, lo
que en la práctica es bastante inespecífico.

\textbf{Después del refinamiento}, los patrones mejoran su fitness y se vuelven más informativos:

\begin{table}[H]
\centering
\footnotesize
\begin{tabular}{P{1cm}P{5.5cm}P{2cm}P{2cm}}
\toprule
ID & Patrón & Fitness & Secuencias \\
\midrule
A1 & \texttt{L-x(1,3)-Y-x-[GST]-x-R}          & 14.1163 & 72 \\
B2 & \texttt{F-x(4,5)-D-x(1,2)-H-x(2)-[ILMV]} & 13.9048 & 70 \\
C3 & \texttt{F-x(4,5)-D-x(0,2)-H-x(2)-[ILMV]} & 13.4048 & 70 \\
D4 & \texttt{Y-x-S-x-R}                         & 12.5102 & 93 \\
\bottomrule
\end{tabular}
\caption{Patrones de PRATT después del refinamiento (Top 4 por fitness).}
\end{table}

El patrón \textbf{B2} (\texttt{F-x(4,5)-D-x(1,2)-H-x(2)-[ILMV]}) es el biológicamente más
interesante: termina en una H que corresponde exactamente a la histidina proximal del hemo
(posición 103 del HMM). En la Proteína X de \textit{C. tetani} el match es
\texttt{FfnslfDi-HvvM} en las posiciones 93--104 de la secuencia, justo alrededor del
residuo His que identificamos como el más conservado del dominio.

Localizando en el logo del HMM: la F del inicio del patrón B2 corresponde a residuos
hidrofóbicos conservados alrededor de la posición 96, la D aparece como una posición
moderadamente conservada en esa región, y la H final es la columna más alta del logo en
posición 103. El refinamiento logró capturar exactamente el bloque funcional del sitio de
unión al hemo.

% ============================================================
\section{Ejercicio 3d}

\begin{consigna}
Busque el patrón de PRATT en PROSITE. ¿Qué tipo de proteínas encuentra? Compare con el patrón
manual.
\end{consigna}

Buscamos el patrón B2 de PRATT (\texttt{F-x(4,5)-D-x(1,2)-H-x(2)-[ILMV]}) en ScanProsite
contra Swiss-Prot (574.627 entradas, release 2026\_01). El resultado alcanzó el límite de
10.000 secuencias con \textbf{10.181 hits en total}. Los primeros resultados incluyen proteínas
completamente no relacionadas con HNOB: transportadores ABC de \textit{Arabidopsis}, lectinas
de arroz (Ricin B-like lectins), oxidasas de hongos, toxinas dérmicas de arañas
(\textit{Loxosceles}), proteínas de virus de insectos (FV-3, ASFV), y receptores de adenosina
humanos, entre otras. No encontramos ninguna ``methyl-accepting chemotaxis protein'' ni proteína
de unión a hemo entre los primeros resultados.

El número esperado de matches al azar es $\sim$2.059 por cada 100.000 secuencias, lo que
es consistente con obtener $\sim$10.000 hits en las 574.627 entradas de Swiss-Prot. Es decir,
prácticamente todos los hits son falsos positivos generados al azar.

Al comparar los tres enfoques:

\begin{itemize}
  \item El \textbf{patrón manual} (\texttt{H-x(11,13)-P-x(14,16)-Y}) también llegó al límite
        de 10.000 hits con proteínas no relacionadas. El problema ahí eran las ventanas de
        $\sim$30 residuos que permiten demasiada variabilidad.

  \item El \textbf{patrón de PRATT refinado} (\texttt{F-x(4,5)-D-x(1,2)-H-x(2)-[ILMV]}) tiene
        ventanas más cortas y un bloque más contiguo, pero igualmente llega a 10.000 hits. Las
        tres posiciones fijas (F, D, H) separadas por 4--5 y 1--2 wildcards son insuficientes
        para discriminar un dominio específico en una base de datos de 574.000 secuencias.
\end{itemize}

Lo que aprendemos de este ejercicio es que ni el patrón manual ni el de PRATT son suficientemente
específicos solos para identificar el dominio HNOB. La diferencia real con el perfil HMM de Pfam
(que sí detecta correctamente el dominio) es que el perfil usa \textbf{todas las posiciones del
alineamiento simultáneamente}, no solo 3 o 4 residuos ancla. Cuando una proteína tiene que cumplir
la distribución de aminoácidos en 180 posiciones al mismo tiempo, los falsos positivos se eliminan.
Los patrones de PROSITE son útiles para sitios activos muy conservados (como el PROSITE de la
serina proteasa: \texttt{[LIVM]-[ST]-A-[STAG]-H-C}), pero para dominios enteros el HMM gana
siempre.
